% GENERATED FILE. Edit canonical lessons, then run npm run generate:books.
% canonical-source-hash: fnv1a64:70c5822a2ea0b658
% canonical-lessons: MW-C80-hamesa, MW-C80-kadi, MW-C80-jaldi, MW-C80-der, MW-C80-paila

\chapter{Always, Sometimes, Quickly, Delay, Before}
\label{ch:mw-always-sometimes-quickly-delay-before}
\hlchaptermodality{80}

\begin{chapteropening}
\textbf{By the end of this chapter:} \emph{I can say always, sometimes, ever, quickly, a delay and before.}

Complete the last lesson of chapter 80: I can say always, sometimes, ever, quickly, a delay and before.
\end{chapteropening}

\section[\texorpdfstring{\mw{हमेसा} (hamesā)}{हमेसा (hamesā)}]{\mw{हमेसा} (hamesā) — always}

\label{lesson:MW-C80-hamesa}



\begin{quote}
Before the new one: say the Marwadi for yesterday; tomorrow, then the Marwadi for now.
\end{quote}

\subsection*{You'll want to know: \mw{हमेसा}}
\textbf{\mw{हमेसा}} — \emph{hamesā} — \textquotedblleft{}always\textquotedblright{}.

\begin{cousinweb}
\textbf{\mw{हमेसा}} (\emph{hamesā}) is a loan, usually traced to Persian \emph{hamēša}, \textquotedblleft{}always\textquotedblright{}.
\end{cousinweb}

\begin{grammarlens}[title={What you've built}]
One more word for this chapter.
\end{grammarlens}

\subsection*{Your turn}
\begin{itemize}
\raggedright
  \item \emph{Say it:} \emph{hamesā}
  \item \emph{Say it:} \emph{hamesā}, once more
  \item \emph{Say it:} say \emph{hamesā}
  \item \emph{Recall:} say the Marwadi for yesterday; tomorrow, then the Marwadi for now, then say \emph{hamesā} again
\end{itemize}

\begin{tcolorbox}[breakable,colback=teal!4,colframe=teal!35!black,title={Before you move on}]
What does \mw{हमेसा} mean? (\textquotedblleft{}Always\textquotedblright{}.) Say it once more.
\end{tcolorbox}
\section[\texorpdfstring{\mw{कदी} (kadī)}{कदी (kadī)}]{\mw{कदी} (kadī) — sometimes, ever}

\label{lesson:MW-C80-kadi}



\begin{quote}
Before the new one: say the Marwadi for now, then the Marwadi for always.
\end{quote}

\subsection*{You'll want to know: \mw{कदी}}
\textbf{\mw{कदी}} — \emph{kadī} — \textquotedblleft{}sometimes, ever\textquotedblright{}.

\begin{cousinweb}
CDIAL has no entry for \textbf{\mw{कदी}} (\emph{kadī}), so its history is not settled. Its \emph{k-} is the old question stem: CDIAL lists Hindi \emph{kabhī}, \textquotedblleft{}ever\textquotedblright{}, under Sanskrit \emph{evam eva}, \textquotedblleft{}exactly so\textquotedblright{}, crossed with the stem \emph{ka-} (CDIAL 2528), but nothing with the shape of \textbf{\mw{कदी}} (\emph{kadī}).
\end{cousinweb}

\begin{grammarlens}[title={What you've built}]
One more word for this chapter.
\end{grammarlens}

\subsection*{Your turn}
\begin{itemize}
\raggedright
  \item \emph{Say it:} \emph{kadī}
  \item \emph{Say it:} \emph{kadī}, once more
  \item \emph{Say it:} say \emph{kadī}
  \item \emph{Recall:} say the Marwadi for now, then the Marwadi for always, then say \emph{kadī} again
\end{itemize}

\begin{tcolorbox}[breakable,colback=teal!4,colframe=teal!35!black,title={Before you move on}]
What does \mw{कदी} mean? (\textquotedblleft{}Sometimes, ever\textquotedblright{}.) Say it once more.
\end{tcolorbox}
\section[\texorpdfstring{\mw{जल्दी} (jaldī)}{जल्दी (jaldī)}]{\mw{जल्दी} (jaldī) — quickly}

\label{lesson:MW-C80-jaldi}



\begin{quote}
Before the new one: say the Marwadi for always, then the Marwadi for sometimes.
\end{quote}

\subsection*{You'll want to know: \mw{जल्दी}}
\textbf{\mw{जल्दी}} — \emph{jaldī} — \textquotedblleft{}quickly\textquotedblright{}.

\begin{cousinweb}
\textbf{\mw{जल्दी}} (\emph{jaldī}) is a loan, usually traced to Persian \emph{jald}, \textquotedblleft{}quick\textquotedblright{}. Persian \emph{jald} is itself of Arabic origin; Marwadi added the ending \emph{-ī}.
\end{cousinweb}

\begin{grammarlens}[title={What you've built}]
One more word for this chapter.
\end{grammarlens}

\subsection*{Your turn}
\begin{itemize}
\raggedright
  \item \emph{Say it:} \emph{jaldī}
  \item \emph{Say it:} \emph{jaldī}, once more
  \item \emph{Say it:} say \emph{jaldī}
  \item \emph{Recall:} say the Marwadi for always, then the Marwadi for sometimes, then say \emph{jaldī} again
\end{itemize}

\begin{tcolorbox}[breakable,colback=teal!4,colframe=teal!35!black,title={Before you move on}]
What does \mw{जल्दी} mean? (\textquotedblleft{}Quickly\textquotedblright{}.) Say it once more.
\end{tcolorbox}
\section[\texorpdfstring{\mw{देर} (der)}{देर (der)}]{\mw{देर} (der) — a delay; late}

\label{lesson:MW-C80-der}



\begin{quote}
Before the new one: say the Marwadi for sometimes, then the Marwadi for quickly.
\end{quote}

\subsection*{You'll want to know: \mw{देर}}
\textbf{\mw{देर}} — \emph{der} — \textquotedblleft{}a delay; late\textquotedblright{}.

\begin{cousinweb}
\textbf{\mw{देर}} (\emph{der}) is a loan, usually traced to Persian \emph{dēr}, \textquotedblleft{}late; a long time\textquotedblright{}.
\end{cousinweb}

\begin{grammarlens}[title={What you've built}]
One more word for this chapter.
\end{grammarlens}

\subsection*{Your turn}
\begin{itemize}
\raggedright
  \item \emph{Say it:} \emph{der}
  \item \emph{Say it:} \emph{der}, once more
  \item \emph{Say it:} say \emph{der}
  \item \emph{Recall:} say the Marwadi for sometimes, then the Marwadi for quickly, then say \emph{der} again
\end{itemize}

\begin{tcolorbox}[breakable,colback=teal!4,colframe=teal!35!black,title={Before you move on}]
What does \mw{देर} mean? (\textquotedblleft{}A delay; late\textquotedblright{}.) Say it once more.
\end{tcolorbox}
\section[\texorpdfstring{\mw{पैलां} (pailā̃)}{पैलां (pailā̃)}]{\mw{पैलां} (pailā̃) — before}

\label{lesson:MW-C80-paila}



\begin{quote}
Before the new one: say the Marwadi for quickly, then the Marwadi for a delay; late.
\end{quote}

\subsection*{You'll want to know: \mw{पैलां}}
\textbf{\mw{पैलां}} — \emph{pailā̃} — \textquotedblleft{}before\textquotedblright{}.

\begin{cousinweb}
\textbf{\mw{पैलां}} (\emph{pailā̃}) is a tadbhava, a word inherited and reshaped by sound change. It comes from \emph{*prathila}, \textquotedblleft{}first\textquotedblright{}, a form CDIAL reconstructs from its descendants. CDIAL lists Old Marwari \emph{pahalo} under it (CDIAL 8652). \textbf{\mw{पैलां}} (\emph{pailā̃}), \textquotedblleft{}before\textquotedblright{}, is a form of the word for \textquotedblleft{}first\textquotedblright{}, built on Sanskrit \emph{prathama}, \textquotedblleft{}first\textquotedblright{}.
\end{cousinweb}

\begin{grammarlens}[title={What you've built}]
One more word for this chapter.
\end{grammarlens}

\subsection*{Your turn}
\begin{itemize}
\raggedright
  \item \emph{Say it:} \emph{pailā̃}
  \item \emph{Say it:} \emph{pailā̃}, once more
  \item \emph{Say it:} say \emph{pailā̃}
  \item \emph{Recall:} say the Marwadi for quickly, then the Marwadi for a delay; late, then say \emph{pailā̃} again
\end{itemize}

\begin{tcolorbox}[breakable,colback=teal!4,colframe=teal!35!black,title={Before you move on}]
What does \mw{पैलां} mean? (\textquotedblleft{}Before\textquotedblright{}.) Say it once more.
\end{tcolorbox}
